% !TEX TS-program = xelatex
% !TEX encoding = UTF-8 Unicode
% !Mode:: "TeX:UTF-8"

\documentclass{resume}
\usepackage{zh_CN-Adobefonts_external} % Simplified Chinese Support using external fonts (./fonts/zh_CN-Adobe/)
%\usepackage{zh_CN-Adobefonts_internal} % Simplified Chinese Support using system fonts
\usepackage{linespacing_fix} % disable extra space before next section
\usepackage{cite}

\begin{document}
\pagenumbering{gobble} % suppress displaying page number

\name{王力}

\basicInfo{
\email{liwang1@link.cuhk.edu.cn} \textperiodcentered\
\faDesktop{\href{https://wwwwwli.github.io/}{~王力个人主页}} \textperiodcentered\
\faGraduationCap {\href{https://scholar.google.com/citations?user=NLQPDrsAAAAJ&hl}{~谷歌学术}} }

\small \textbf{研究方向：} Safety, Audio Deepfake Detection, Audio Language Model
 
\section{教育经历}
\datedsubsection{\textbf{香港中文大学（深圳）}}{2023年9月 -- 至今}
\begin{itemize}
\item 研究领域：音频深度伪造检测与反欺诈
\item 导师：\href{https://sds.cuhk.edu.cn/teacher/641}{武执政} 和 \href{https://sds.cuhk.edu.cn/teacher/498}{李海洲}
\end{itemize}
\datedsubsection{\textbf{北京大学}（免试硕士研究生）}{2019年9月 -- 2022年6月}
\begin{itemize}
\item 计算机科学与技术硕士研究生
\item 研究领域：口语关键词识别和说话人验证
\item 导师：\href{https://web.pkusz.edu.cn/adsp/}{邹月娴}
\end{itemize}
\datedsubsection{\textbf{电子科技大学}}{2015年9月 -- 2019年6月}
\begin{itemize}
\item 信息与软件工程学士
\item GPA：3.69/4.00，排名：8/162（前4.9\%）
\item 主修课程：随机信号分析（99分），数字信号处理（95分）
\end{itemize}

\section{发表论文}
\begin{enumerate}\itemsep 0.5em
  \item \textbf{Li Wang}, Xiao Lei, Haorui He, Lei Wang, Jie Shi, Zhizheng Wu. \href{https://ieeexplore.ieee.org/document/11334038}{[PDF]} \\Over-the-Air Adversarial Attacks and Detection for Automatic Speaker Verification. \textit{IEEE/ACM Transactions on Audio, Speech, and Language Processing (TASLP), 2026.}
  \item Peizhuo Liu*, \textbf{Li Wang*}, Renqiang He*, Haorui He, Lei Wang, Huadi Zheng, Jie Shi, Tong Xiao, Zhizheng Wu. \href{https://arxiv.org/abs/2409.11308}{[PDF]} \\SpMis: An Investigation of Synthetic Spoken Misinformation Detection. \textit{Proceedings of the IEEE Spoken Language Technology Workshop (SLT 2024).} (\textbf{Best Paper Finalist, Top 2.5\%})
  \item \textbf{Li Wang}, Jiaqi Li, Yuhao Luo, Jiahao Zheng, Lei Wang, Hao Li, Ke Xu, Chengfang Fang, Jie Shi, Zhizheng Wu. \href{https://arxiv.org/abs/2310.05369}{[PDF]} \\AdvSV: An Over-the-Air Adversarial Attack Dataset for Speaker Verification. \textit{Proceedings of the 49th IEEE International Conference on Acoustics, Speech and Signal Processing (ICASSP 2024).}
 \item Jiaqi Li, \textbf{Li Wang}, Liumeng Xue, Lei Wang, Zhizheng Wu. \href{https://arxiv.org/abs/2310.05354}{[PDF]} \\An Initial Investigation of Neural Replay Simulator for Over-the-Air Adversarial Perturbations to Automatic Speaker Verification. \textit{Proceedings of the 49th IEEE International Conference on Acoustics, Speech and Signal Processing (ICASSP 2024).}
  \item \textbf{Li Wang}, Rongzhi Gu, Weiji Zhuang, Peng Gao, Yujun Wang, Yuexian Zou. \href{https://ieeexplore.ieee.org/document/9747878}{[PDF]} \\Learning Decoupling Features Through Orthogonality Regularization. \textit{Proceedings of the 47th IEEE International Conference on Acoustics, Speech and Signal Processing (ICASSP 2022).}
  \item \textbf{Li Wang}, Rongzhi Gu, Nuo Chen, Yuexian Zou. \href{https://www.isca-speech.org/archive/interspeech_2021/wang21da_interspeech.html}{[PDF]} \\ Text Anchor Based Metric Learning for Small-Footprint Keyword Spotting. \textit{Proceedings of the 22nd Annual Conference of the International Speech Communication Association (Interspeech 2021).}
\end{enumerate}

\section{工作经历}
\datedsubsection{\textbf{深圳市大数据研究院（SRIBD）}}{2022年11月 -- 2023年8月}
{\small \role{研究助理，大数据基础理论与算法研究所，中国深圳}{}}
负责与华为合作的演讲者验证技术中物理攻击的检测项目
\begin{itemize}
\item 生成110万条面向说话人验证的物理攻击数据
\item 确定了防御方法的技术路线，得到了华为的充分肯定和支持
\item 为说话人验证技术的物理攻击开发了可靠的防御机制
\item 在整个项目过程中展现出了出色的领导力、问题解决能力和技术能力
\end{itemize}

\datedsubsection{\textbf{粤港澳大湾区数字经济研究院（IDEA）}}{2022年7月 -- 2022年10月}
{\small \role{金融机器学习工程师，AI金融与深度学习研究中心，中国深圳}{}}
创建金融因子（特征）并设计深度学习模型以预测股票回报

\section{实习经历}
\datedsubsection{\textbf{字节跳动AILab}}{2021年7月 -- 2022年3月}
\begin{itemize}
\item 课题：音频指纹和音乐分割
\item 工作描述：提出了音频指纹的时间定位算法，将准确率从95.80\%提高到97.90\%，时间定位的IOU为0.62
\end{itemize}

\datedsubsection{\textbf{华为MTI音频工程部}}{2021年4月 -- 2022年6月}
\begin{itemize}
\item 课题：音乐分离
\item 工作描述：复现D3Net，使用MUSDB数据集进行训练和测试，将人声、鼓声、低音和其他音轨分离，测试平均SDR为6.43
\end{itemize}

\datedsubsection{\textbf{腾讯AILab}}{2020年5月 -- 2020年11月}
\begin{itemize}
\item 课题：语音分离和说话人验证
\item 工作描述：在说话人验证任务中测试语音分离模型GALR的性能，整理了200万条音频数据的说话人信息
\end{itemize}

\section{项目经历}
\datedsubsection{\textbf{人工智能语音唤醒技术研究}}{2021年1月 -- 2022年1月}
\role{负责人}{研究生期间与小米公司合作的协作项目}
\begin{itemize}
\item 与小米公司沟通，跟踪项目进展，了解数据需求
\item 优化网络设计和数据增强，提高模型性能
\item 将误拒绝率从20\%降低到2\%，准确率达到99.7\%
\item 参与个性化唤醒任务的合作研究，并作为第一作者发表ICASSP2022论文
\end{itemize}

% \datedsubsection{\textbf{疲劳驾驶智能监测系统}}{2017年9月 -- 2018年6月}
% \role{负责人}{小组共5名成员}
% \begin{itemize}
% \item 开发疲劳驾驶监测系统，并将其部署到Raspberry Pi开发板上
% \item 白天和夜间的疲劳检测准确率分别达到73\%和82\%
% \item 开发了一个安卓应用程序，实时提醒驾驶员疲劳状态并提供警示提醒
% \item 该项目在全国大学生信息安全竞赛中获得全国二等奖
% \end{itemize}

\section{服务经历}
\datedsubsection{\textbf{审稿人}}{}
\begin{itemize}
  \item IEEE/ACM Transactions on Audio, Speech, and Language Processing (IF 4.1)
  \item IEEE International Conference on Acoustics, Speech and Signal Processing
  \item IEEE Signal Processing Letters (IF 3.201)
  \item 全国人机语音通讯学术会议（NCMMSC）
  \item International Conference on Asian Language Processing (IALP)
  \item Computer Speech \& Language (IF 3.252)
  \item Workshop on Automatic speech Recognition and Understanding
\end{itemize}

\datedsubsection{\textbf{标准化工作}}{}
\begin{itemize}
  \item 参与全国网络安全标准化技术委员会（TC260）《人工智能生成合成内容标识方法 文件元数据隐式标识 文本文件》的制定工作。
\end{itemize}

\datedsubsection{\textbf{学生志愿者}}{}
\begin{itemize}
  \item IEEE Spoken Language Technology Workshop 2024
\end{itemize}

\datedsubsection{\textbf{助教}}{}
\begin{itemize}
  \item 2023 秋季, CSC1001 计算机科学导论：编程方法学, 香港中文大学（深圳）.
  \item 2024 春季, CSC3160/AIR6063 语音与语言处理基础 / 口语语言处理, 香港中文大学（深圳）.
  \item 2024 Fall, DDA3020 机器学习, 香港中文大学（深圳）.
  \item 2025 Spring, DDA3020 机器学习, 香港中文大学（深圳）.

\end{itemize}

\section{荣誉和奖项}
\begin{enumerate}
\item 全国大学生信息安全竞赛国家级二等奖（2018年）
\item 本科生人民一等奖学金（2015年）
\item 本科生人民一等奖学金（2016年）
\item 本科生人民一等奖学金（2017年）
\item 中华人民共和国第九届残疾人运动会优秀志愿者（2015年）
\item 北京大学校运会4x100米亚军（2021年）
\end{enumerate}

\section{技能}
\begin{itemize}
\item 二胡十级
\item 编程语言：熟练使用Python，熟悉Java，了解C++
\item 深度学习框架：熟练使用PyTorch
\end{itemize}

%% Reference
%\newpage
%\bibliographystyle{IEEETran}
%\bibliography{mycite}
\end{document}
